\documentclass[12pt, a4paper]{article}

% --- Pacotes Fundamentais ----
\usepackage[utf8]{inputenc}
\usepackage[T1]{fontenc}
\usepackage[brazil]{babel}
\usepackage{geometry}
\usepackage{titlesec}
\usepackage{fancyhdr}
\usepackage{graphicx}
\usepackage{xcolor}
\usepackage{booktabs}
\usepackage{longtable}
\usepackage{tabularx}
\usepackage{colortbl}
\usepackage{float}
\usepackage{parskip}
\usepackage{lmodern}
\usepackage{lastpage}
\usepackage{tikz}
\usepackage{pgfplots}
\usepackage{amssymb}
\usepackage{needspace}
\usepackage{hyperref}
\usepackage{enumitem}
\pgfplotsset{compat=1.18}

% Configuração de hyperlinks
\hypersetup{
    colorlinks=true,
    linkcolor=primaryColor,
    urlcolor=accentColor,
    pdftitle={Processo Completo: Teste de Criativos, Ofertas e Funis Low Ticket},
    pdfauthor={Operação Low Ticket}
}

% --- Configurações de Layout ---
\geometry{left=2.5cm, right=2.5cm, top=4.5cm, bottom=2cm, headsep=1.5cm, headheight=60pt}

% --- Definição de Cores ---
\definecolor{primaryColor}{RGB}{0, 51, 102}   % Azul Marinho
\definecolor{secondaryText}{RGB}{80, 80, 80}  % Cinza Escuro
\definecolor{accentColor}{RGB}{178, 34, 34}   % Vermelho Tijolo
\definecolor{lightGray}{RGB}{245, 245, 245}
\definecolor{tableHeader}{RGB}{230, 230, 250}
\definecolor{successGreen}{RGB}{34, 139, 34}  % Verde para sucesso
\definecolor{warningOrange}{RGB}{255, 140, 0} % Laranja para atenção

% --- Configuração de Cabeçalho e Rodapé ---
\pagestyle{fancy}
\fancyhf{}

% Cabeçalho
\fancyhead[C]{%
    \begin{tabularx}{\textwidth}{@{} X r @{}}
        \textcolor{primaryColor}{\textbf{\footnotesize PROCESSO COMPLETO: CRIATIVOS, OFERTAS E FUNIS LOW TICKET}} & \scriptsize \textit{Guia Operacional} \\
        \textcolor{secondaryText}{\scriptsize Teste $\rightarrow$ Validação $\rightarrow$ Escala} & {\scriptsize \textbf{Data Base: Março/2026}}
    \end{tabularx}%
}

% Rodapé
\fancyfoot[L]{\scriptsize \textcolor{secondaryText}{Referência: \textbf{Victor Xisto} | Operação Low Ticket}}
\fancyfoot[R]{\normalsize Página \textbf{\thepage} de \textbf{\pageref{LastPage}}}

% --- Linha separadora do CABEÇALHO ---
\renewcommand{\headrulewidth}{1.5pt}
\renewcommand{\headrule}{%
    \vspace{0.1cm}
    \hbox to\headwidth{\color{primaryColor}\leaders\hrule height \headrulewidth\hfill}%
}

% --- Linha separadora do RODAPÉ ---
\renewcommand{\footrulewidth}{1.5pt}
\renewcommand{\footrule}{%
    \color{primaryColor}%
    \hbox to\headwidth{\leaders\hrule height \footrulewidth\hfill}%
    \vspace{0.2cm}
}

% --- Formatação de Títulos ---
\titleformat{\section}
{\color{primaryColor}\normalfont\Large\bfseries}
{\thesection}{1em}{}[{\titlerule[0.8pt]}]

\titleformat{\subsection}
{\color{primaryColor}\normalfont\large\bfseries}
{\thesubsection}{1em}{}

% --- Macro para Moeda ---
\newcommand{\reais}[1]{R\$ #1}

% --- INÍCIO DO DOCUMENTO ---
\begin{document}

% --- CAPA ---
\begin{titlepage}
    \centering
    \vspace*{2.5cm}
    {\Huge \textbf{\textcolor{primaryColor}{PROCESSO COMPLETO}}} \\
    \vspace{0.5cm}
    {\Large \textbf{TESTE DE CRIATIVOS, OFERTAS E FUNIS LOW TICKET}} \\
    \vspace{0.3cm}
    {\large Da concepção da oferta à escala com lucro no front-end} \\
    \vspace{1.5cm}
    \rule{10cm}{1.5pt} \\
    \vspace{0.5cm}
    \colorbox{lightGray}{%
        \parbox{0.9\textwidth}{%
            \centering
            \vspace{0.3cm}
            {\Large \textbf{GUIA OPERACIONAL PARA APLICAÇÃO IMEDIATA}} \\
            \vspace{0.3cm}
            {\large Mentalizar o processo, validar ofertas e escalar com previsibilidade} \\
            \vspace{0.3cm}
        }%
    }
    \vspace{1.5cm}

    \begin{flushleft}
    \textbf{PILARES DO PROCESSO:} \\
    \vspace{0.2cm}
    \textcolor{primaryColor}{\textbf{1}} - Concepção e Modelagem de Ofertas (Quiz + App) \\
    \textcolor{primaryColor}{\textbf{2}} - Criação e Teste de Criativos (Orgânico + TikTok) \\
    \textcolor{primaryColor}{\textbf{3}} - Estrutura de Tráfego (ABO/CBO + Teste de Orçamento) \\
    \textcolor{primaryColor}{\textbf{4}} - Validação e Escala (Conta, Estrutura, Budget)
    \end{flushleft}

    \vfill
    {\small Documento operacional compilado a partir de metodologia validada em operações reais de low ticket, com foco em funis de quiz, aplicativos como entregável, criativos estilo orgânico/TikTok e estratégias de tráfego pago com taxa de assertividade de 80\% em novas ofertas.} \\
    \vspace{1.5cm}

    % Box de Créditos
    \colorbox{accentColor}{%
        \parbox{0.9\textwidth}{%
            \centering
            \vspace{0.6cm}
            \textcolor{white}{\Large \textbf{REFERÊNCIA: VICTOR XISTO}} \\
            \vspace{0.2cm}
            \textcolor{white}{\normalsize Operação Low Ticket | 140K+ alunos | Escala 100K+/dia} \\
            \vspace{0.6cm}
        }%
    }

    \vspace{0.2cm}
    \textbf{Data Base:} Março de 2026
\end{titlepage}

% --- CONTEÚDO ---
\section{Visão de Mercado: Por Que Low Ticket com Recorrência}

O modelo mais valorizado do mercado digital atual combina \textbf{low ticket + recorrência + aplicativo}. Operações neste formato possuem valuation até \textbf{7 vezes maior} que modelos tradicionais de infoproduto.

\subsection{Por Que Aplicativos e Recorrência}

\begin{table}[H]
    \centering
    \renewcommand{\arraystretch}{1.4}
    \begin{tabularx}{\textwidth}{X X}
        \toprule
        \rowcolor{tableHeader} \textbf{Modelo Tradicional} & \textbf{Modelo App + Recorrência} \\
        \midrule
        Venda única (ticket R\$27-R\$97) & Assinatura mensal + upsells internos \\
        \rowcolor{lightGray}
        LTV limitado ao front-end & LTV alto (meses de recorrência) \\
        Sem equity real & Valuation 7x superior \\
        \rowcolor{lightGray}
        Entregável genérico (e-book/PDF) & Aplicativo personalizado com IA \\
        Difícil de vender a operação & Propostas reais de aquisição \\
        \bottomrule
    \end{tabularx}
\end{table}

\subsection{A Oportunidade com IA}

Antigamente, criar um aplicativo exigia equipe de programação e meses de desenvolvimento. Hoje, com ferramentas como \textbf{Replit} e \textbf{Lovable}, é possível criar um MVP funcional em \textbf{2 a 3 dias}.

\begin{quote}
\textcolor{primaryColor}{\textbf{``Toda vez que a gente cria um produto novo, a gente já cria na intenção dele virar um app, um negócio que tem recorrência.''}}
\end{quote}

\textcolor{accentColor}{\textbf{Modelo de Holding:}} Crie múltiplos produtos do mesmo nicho sob a mesma empresa. Exemplo: uma holding de saúde com 3 ofertas (emagrecimento feminino, masculino e saúde geral), acumulando público e alunos. Um dia essa empresa pode ser vendida com valuation significativo.

\needspace{8cm}

\section{O Funil de Quiz: A Estrutura que Converte}

O quiz é o formato de funil mais validado para low ticket. A regra fundamental é: \textbf{80/20 é personalização}.

\subsection{Estrutura do Quiz em 3 Blocos}

\begin{table}[H]
    \centering
    \renewcommand{\arraystretch}{1.5}
    \begin{tabularx}{\textwidth}{c X X}
        \toprule
        \rowcolor{tableHeader} \textbf{Bloco} & \textbf{Objetivo} & \textbf{Conteúdo} \\
        \midrule
        \textbf{Início} & Pesquisa / Coleta de dados & Perguntas sobre perfil, objetivos, situação atual do lead \\
        \rowcolor{lightGray}
        \textbf{Meio} & Projeção de resultado & Mostrar como o lead vai ficar, resultados possíveis, antes e depois \\
        \textbf{Final} & Entrega personalizada + Oferta & Plano/resultado com base nas respostas + preço baixo por impulso \\
        \bottomrule
    \end{tabularx}
\end{table}

\subsection{Por Que o Quiz Funciona para Low Ticket}

\begin{itemize}
    \item \textbf{O lead se enxerga no futuro:} A personalização faz ele visualizar o resultado antes de comprar
    \item \textbf{Percepção de valor alto:} O quiz entrega algo que parece valer muito mais do que o preço cobrado
    \item \textbf{Compra por impulso:} No final do quiz, o lead já está emocionalmente investido e o preço é baixo
    \item \textbf{Volume:} O quiz permite trazer muito tráfego barato (``teste grátis'') e converter no funil
\end{itemize}

\textcolor{accentColor}{\textbf{Regra de Ouro:}} O quiz de personalização supera consistentemente quizzes baseados apenas em depoimentos ou mini-VSL. Sempre maximize o benefício com base no que a pessoa respondeu.

\needspace{10cm}

\section{Concepção e Modelagem de Ofertas}

A tomada de decisão sobre qual oferta criar segue dois eixos: \textbf{timing (sazonalidade)} e \textbf{modelagem de ofertas validadas}.

\subsection{Eixo 1: Timing e Sazonalidade}

Ofertas sazonais são ouro para low ticket porque você entra no leilão barato --- poucos anunciantes competem fora da temporada.

\begin{table}[H]
    \centering
    \renewcommand{\arraystretch}{1.4}
    \begin{tabularx}{\textwidth}{X X X}
        \toprule
        \rowcolor{tableHeader} \textbf{Evento Sazonal} & \textbf{Tipo de Produto} & \textbf{Vantagem} \\
        \midrule
        ENEM / Vestibular & App com resumos + comunidade & Vende todo ano, acumula público \\
        \rowcolor{lightGray}
        Páscoa & Receitas de ovos / renda extra & Leilão barato, demanda garantida \\
        Natal / Fim de Ano & Renda extra, presentes DIY & Pico de consumo \\
        \rowcolor{lightGray}
        Verão / Janeiro & Emagrecimento, treinos & Demanda sazonal forte \\
        \bottomrule
    \end{tabularx}
\end{table}

\textcolor{primaryColor}{\textbf{Conceito de ``Sazonal Infinito'':}} Um produto de ENEM, por exemplo, vende durante a temporada, acumula seguidores e alunos, e no ano seguinte vende mais --- com app mais sofisticado e base de público maior. Suba antes da concorrência, com oferta validada e produto bem feito.

\subsection{Eixo 2: Modelagem de Ofertas Validadas}

O processo não é copiar --- é \textbf{entender princípios e adaptar}:

\begin{enumerate}
    \item \textbf{Pesquise mercados externos:} Faça spy em outros países (EUA, Latam) e identifique ofertas escalando
    \item \textbf{Identifique o mecanismo validado:} Qual o gatilho, a promessa, a estrutura que funciona
    \item \textbf{Transforme o entregável:} Converta de produto físico/VSL/Nutra para \textbf{aplicativo low ticket}
    \item \textbf{Mude o Nome Chiclete:} Crie um nome novo que pareça inovador (mesmo mecanismo por trás)
    \item \textbf{Monte o quiz de personalização:} Adapte o funil para o formato quiz com a estrutura de 3 blocos
\end{enumerate}

\needspace{8cm}

\subsection{Conversão de VSL/Nutra para Low Ticket}

Este é o processo de maior valor: pegar uma oferta validada em outro formato e transformá-la em low ticket.

\begin{table}[H]
    \centering
    \renewcommand{\arraystretch}{1.5}
    \begin{tabularx}{\textwidth}{c X X}
        \toprule
        \rowcolor{tableHeader} \textbf{Etapa} & \textbf{Oferta Original (VSL/Nutra)} & \textbf{Adaptação Low Ticket} \\
        \midrule
        \textbf{1} & Mecanismo de conversão (ex: Pink Salt, diabetes) & Manter o mecanismo, mudar o nome chiclete \\
        \rowcolor{lightGray}
        \textbf{2} & Página de VSL longa (15-40 min) & Quiz de personalização (2-4 min) \\
        \textbf{3} & Produto físico (cápsula, suplemento) & Aplicativo/plano digital personalizado \\
        \rowcolor{lightGray}
        \textbf{4} & Ticket alto (R\$150-R\$400) & Ticket baixo (R\$17-R\$47) + recorrência \\
        \textbf{5} & Funil de consciência longo & Funil de impulso curto (quiz $\rightarrow$ oferta) \\
        \bottomrule
    \end{tabularx}
\end{table}

\textcolor{accentColor}{\textbf{Atenção:}} Não basta trocar o formato --- o entregável precisa ser bom. Não adianta pegar uma oferta de Nutra e entregar um e-book. O formato validado é \textbf{aplicativo, plano personalizado ou desafio}.

\subsection{O Poder do Nome Chiclete}

O nome chiclete é 80/20 na percepção de novidade da oferta:

\begin{itemize}
    \item Mesmo mecanismo validado + nome chiclete novo = \textbf{oferta que parece inovadora}
    \item Múltiplos nomes chiclete validados permitem atingir \textbf{fatias de público diferentes}
    \item Monte um banco de \textbf{10 nomes chiclete} e teste sistematicamente
    \item O nome chiclete no criativo + no quiz = melhor performance
\end{itemize}

\begin{quote}
\textcolor{primaryColor}{\textbf{``Por que a gente consegue escalar 100K/dia? Porque a gente tem vários nomes chiclete validados que pegam fatias de público diferentes. A pessoa vê o quiz de diferentes anunciantes e nem percebe que é o mesmo quiz.''}}
\end{quote}

\needspace{10cm}

\section{Criativo para Low Ticket: O Formato que Escala}

O criativo é o 80/20 da operação. Após ter oferta validada, o jogo é \textbf{criativo e tráfego}.

\subsection{Diferenças Fundamentais: Low Ticket vs. Nutra/VSL}

\begin{table}[H]
    \centering
    \renewcommand{\arraystretch}{1.4}
    \begin{tabularx}{\textwidth}{X X}
        \toprule
        \rowcolor{tableHeader} \textbf{Criativo Nutra/VSL} & \textbf{Criativo Low Ticket} \\
        \midrule
        Longo (1-3 min ou mais) & Curto (até 50 segundos) \\
        \rowcolor{lightGray}
        Conscientiza e prova no criativo & Gera curiosidade e joga pro quiz \\
        Mostra produto físico & Mostra a pessoa usando o app \\
        \rowcolor{lightGray}
        Formato mais produzido & Formato 100\% orgânico (estilo TikTok) \\
        Objetivo: vender no criativo & Objetivo: gerar clique para o quiz \\
        \bottomrule
    \end{tabularx}
\end{table}

\subsection{Anatomia do Criativo Low Ticket Validado}

\begin{table}[H]
    \centering
    \renewcommand{\arraystretch}{1.5}
    \begin{tabularx}{\textwidth}{c X c}
        \toprule
        \rowcolor{tableHeader} \textbf{Elemento} & \textbf{Descrição} & \textbf{Duração} \\
        \midrule
        \textbf{Hook} & Gancho viral do TikTok + antes/depois explícito & 0-5s \\
        \rowcolor{lightGray}
        \textbf{Body} & UGC natural: pessoa usando o app / fazendo ação + antes/depois no meio & 5-35s \\
        \textbf{CTA} & ``Teste grátis'' / ``Preencha o quiz e ganhe um presente'' + antes/depois final & 35-50s \\
        \bottomrule
    \end{tabularx}
\end{table}

\subsection{A Técnica do Antes e Depois Repetido}

Essa técnica é o que faz o criativo low ticket converter em volume:

\begin{enumerate}
    \item Coloque um \textbf{antes e depois no início} (logo nos primeiros segundos)
    \item Coloque outro \textbf{antes e depois no meio} (após mostrar o body)
    \item Coloque um terceiro \textbf{antes e depois no final} (antes do CTA)
\end{enumerate}

\textcolor{accentColor}{\textbf{Regra Crítica:}} O antes e depois deve ser \textbf{da mesma pessoa} em todas as 3 aparições. Erro fatal: usar uma mulher loira no início e uma morena no final --- o lead perde a conexão e não converte.

\begin{quote}
\textcolor{primaryColor}{\textbf{``O antes e depois não é um atrás do outro em sequência longa. É um antes e depois que fica aparecendo toda hora em duração curta. Gruda na mente da pessoa --- ela viu 3 provas.''}}
\end{quote}

\needspace{8cm}

\subsection{Regras de Produção de Criativos}

\begin{table}[H]
    \centering
    \renewcommand{\arraystretch}{1.4}
    \begin{tabularx}{\textwidth}{X X}
        \toprule
        \rowcolor{tableHeader} \textbf{Elemento} & \textbf{Padrão Validado} \\
        \midrule
        Duração máxima & 50 segundos \\
        \rowcolor{lightGray}
        Estilo visual & 100\% orgânico / UGC (User Generated Content) / TikTok \\
        Cortes / Transições & A cada 3 segundos \\
        \rowcolor{lightGray}
        Música & Viral do momento (trending) \\
        Hooks & Banco de dados com ganchos do TikTok \\
        \rowcolor{lightGray}
        Antes e Depois & 3x por criativo (início, meio, fim) \\
        Demonstração do App & Pessoa usando fisicamente (não tela estática) \\
        \rowcolor{lightGray}
        CTA principal & ``Teste grátis'' / ``Quiz gratuito'' / ``Presente no final'' \\
        \bottomrule
    \end{tabularx}
\end{table}

\subsection{Processo de Criação de Hooks}

\begin{enumerate}
    \item Mantenha \textbf{2 pessoas dedicadas} a carregar o banco de dados de ganchos do TikTok
    \item Identifique hooks virais e \textbf{replique o gancho} em um body que já deu certo
    \item Combine hooks novos + body validado = criativos com alta taxa de acerto
    \item Exemplo de hook: \textit{``Oi meninas, hoje vocês vão ver como eu sequei minha barriga fazendo esse treino de 7 dias em casa''}
\end{enumerate}

\textcolor{warningOrange}{\textbf{Sobre o CTR:}} Quando o CTR é alto, as métricas secundárias melhoram automaticamente --- CPC cai, CPM cai. O Facebook entrega mais porque vê engajamento. Criativos que mencionam ``teste grátis'' ou ``presente no final'' aumentam significativamente o CTR.

\needspace{10cm}

\section{Estrutura de Tráfego: O Tripé de Testes}

A estratégia de tráfego não é ``usar a estrutura que funciona pro mercado''. Cada conta de anúncio performa de um jeito diferente. O processo exige \textbf{3 testes obrigatórios} em toda nova oferta.

\subsection{Os 3 Testes Obrigatórios}

\noindent
\begin{center}
\resizebox{\textwidth}{!}{%
\begin{tikzpicture}[
    box/.style={
        rounded corners=8pt,
        minimum width=5cm,
        minimum height=4cm,
        inner sep=6pt,
    },
    boxtitle/.style={
        font=\Large\bfseries,
    },
    boxsubtitle/.style={
        font=\large\bfseries,
    },
    boxcontent/.style={
        font=\normalsize,
        text width=4.2cm,
        align=center,
    },
    arrow/.style={
        ->,
        line width=2pt,
        primaryColor,
        >=stealth,
        shorten <=3pt,
        shorten >=3pt,
    },
]
    % Teste 1
    \node[box, fill=primaryColor!12, draw=primaryColor!50, line width=1pt] (t1) at (0,0) {};
    \node[boxtitle, text=primaryColor] at (0, 1.2) {TESTE 1};
    \node[boxsubtitle] at (0, 0.4) {Estrutura};
    \draw[primaryColor!40, line width=0.5pt] (-1.8, 0) -- (1.8, 0);
    \node[boxcontent] at (0, -0.8) {ABO 1-1-1\\[3pt]ABO 1-3-1\\[3pt]CBO 1-3-3};

    % Teste 2
    \node[box, fill=accentColor!10, draw=accentColor!50, line width=1pt] (t2) at (7,0) {};
    \node[boxtitle, text=accentColor] at (7, 1.2) {TESTE 2};
    \node[boxsubtitle] at (7, 0.4) {Orçamento};
    \draw[accentColor!40, line width=0.5pt] (5.2, 0) -- (8.8, 0);
    \node[boxcontent] at (7, -0.8) {R\$17 (baixo)\\[3pt]Valor do ticket\\[3pt]3x o ticket};

    % Teste 3
    \node[box, fill=successGreen!10, draw=successGreen!50, line width=1pt] (t3) at (14,0) {};
    \node[boxtitle, text=successGreen] at (14, 1.2) {TESTE 3};
    \node[boxsubtitle] at (14, 0.4) {Conta de Anúncio};
    \draw[successGreen!40, line width=0.5pt] (12.2, 0) -- (15.8, 0);
    \node[boxcontent] at (14, -0.8) {Conta A\\[3pt]Conta B\\[3pt]\textit{(mesmo criativo)}};

    % Setas
    \draw[arrow] (t1.east) -- (t2.west);
    \draw[arrow] (t2.east) -- (t3.west);
\end{tikzpicture}%
}
\end{center}

\vspace{0.3cm}

\textcolor{accentColor}{\textbf{Por que isso importa:}} Existem operações que foram salvas apenas com mudança de estrutura. Mesmo criativo, mesma copy, mesmo título --- trocou de CBO para ABO (ou vice-versa) e começou a vender. Cada conta performa diferente.

\subsection{Teste de Estrutura}

\begin{table}[H]
    \centering
    \renewcommand{\arraystretch}{1.4}
    \begin{tabularx}{\textwidth}{X X c}
        \toprule
        \rowcolor{tableHeader} \textbf{Estrutura} & \textbf{Quando Usar} & \textbf{Prioridade} \\
        \midrule
        \textbf{ABO 1-1-1} & Teste inicial com orçamento baixo & Alta \\
        \rowcolor{lightGray}
        \textbf{ABO 1-3-1} & Validação de criativo (3 fatias de público) & Alta \\
        \textbf{ABO Multi-Criativo} & 3-4 criativos validados agrupados & Média \\
        \rowcolor{lightGray}
        \textbf{CBO 1-3-1} & Escala com criativo de constância validada & Alta \\
        \textbf{CBO 1-3-3} & Teste de múltiplos criativos em escala & Média \\
        \bottomrule
    \end{tabularx}
\end{table}

\needspace{8cm}

\subsection{Teste de Orçamento: A Estratégia do R\$17}

Esta é a estratégia mais contraintuitiva e mais poderosa validada na operação:

\begin{table}[H]
    \centering
    \renewcommand{\arraystretch}{1.5}
    \begin{tabularx}{\textwidth}{c X X}
        \toprule
        \rowcolor{tableHeader} \textbf{Orçamento} & \textbf{Lógica} & \textbf{Resultado Esperado} \\
        \midrule
        \textcolor{successGreen}{\textbf{R\$17/conjunto}} & Facebook tem o dia todo para gastar pouco. Entrega vendas com CPA mínimo & CPA absurdamente baixo, ROI de 5-7x por conjunto \\
        \rowcolor{lightGray}
        \textbf{Valor do ticket} & Orçamento padrão, permite validação equilibrada & CPA moderado, ROI médio \\
        \textbf{3x o ticket} & Orçamento de escala, mais verba por conjunto & CPA mais alto, volume maior \\
        \bottomrule
    \end{tabularx}
\end{table}

\begin{quote}
\textcolor{primaryColor}{\textbf{``Com R\$17 por conjunto, o Facebook não pode desperdiçar. Ele vai gastar esses R\$17 da forma mais inteligente possível. Ele precisa entregar venda porque precisa que o cliente continue investindo.''}}
\end{quote}

\subsection{Como Aplicar a Estratégia R\$17}

\begin{enumerate}
    \item Pegue um \textbf{criativo validado} (que já vendeu em outra estrutura)
    \item Crie múltiplas campanhas ABO 3:1 com orçamento de R\$17 por conjunto
    \item \textbf{Duplique dezenas de campanhas} com essa estrutura
    \item Resultado: 90\% dos conjuntos vendem, com CPA muito inferior ao da escala com orçamento alto
    \item O faturamento total se iguala ou supera o modelo de ``poucos conjuntos com muito orçamento''
\end{enumerate}

\textcolor{warningOrange}{\textbf{Pré-requisito:}} Funciona melhor com pixel já aquecido. Para quem já roda uma oferta, pegar o criativo campeão e aplicar essa estrutura é o teste mais rápido de implementar.

\subsection{Teste de Conta de Anúncio}

\begin{itemize}
    \item Sempre teste a \textbf{mesma oferta em pelo menos 2 contas} de anúncio diferentes
    \item Existe conta que performa melhor ABO e conta que performa melhor CBO
    \item Contas novas tendem a entregar melhor no início (Facebook quer fidelizar o anunciante)
    \item Se a oferta não valida em uma conta, \textbf{teste em outra} antes de descartar
\end{itemize}

\needspace{8cm}

\section{Fluxo Completo: Da Ideia à Escala}

Este é o mapa mental completo do processo, da concepção da oferta até a escala consolidada.

\subsection{Fase 1: Concepção da Oferta (2-3 dias)}

\begin{enumerate}
    \item Escolha o eixo: \textbf{Timing} (sazonal) ou \textbf{Modelagem} (spy de mercado)
    \item Identifique o mecanismo validado (VSL, Nutra, oferta gringa)
    \item Transforme o entregável em \textbf{aplicativo} (Replit/Lovable, 2-3 dias)
    \item Crie o \textbf{nome chiclete} (banco de 10 nomes, testar progressivamente)
    \item Monte o \textbf{quiz de personalização} (3 blocos: pesquisa $\rightarrow$ projeção $\rightarrow$ oferta)
\end{enumerate}

\subsection{Fase 2: Produção de Criativos (1-2 dias)}

\begin{enumerate}
    \item Selecione 5-10 hooks do banco de dados TikTok
    \item Produza 3-5 criativos de até 50 segundos (orgânico/UGC)
    \item Aplique a técnica do antes/depois repetido (3x, mesma pessoa)
    \item Demonstre o app sendo usado fisicamente (não tela estática)
    \item CTA: ``teste grátis'' / ``quiz gratuito'' / ``presente no final''
\end{enumerate}

\subsection{Fase 3: Testes de Tráfego (3-5 dias)}

\begin{table}[H]
    \centering
    \small
    \renewcommand{\arraystretch}{1.3}
    \begin{tabularx}{\textwidth}{c X X c}
        \toprule
        \rowcolor{tableHeader} \textbf{Dia} & \textbf{Ação} & \textbf{O Que Observar} & \textbf{Decisão} \\
        \midrule
        \textbf{Sex} & Subir criativos (ABO + CBO, 2 contas, 3 orçamentos) & --- & Deixar rodar \\
        \rowcolor{lightGray}
        \textbf{Sáb} & Monitorar métricas iniciais & CTR, CPM, primeiras vendas & Não mexer \\
        \textbf{Dom} & Continuação do teste & ROAS parcial, CPA por conjunto & Não mexer \\
        \rowcolor{lightGray}
        \textbf{Seg} & Primeira análise completa & ROAS acumulado, métricas 2rias & Pausar perdedores \\
        \textbf{Ter-Qua} & Confirmação de constância & ROAS estável? CPA sustentável? & Escalar vencedores \\
        \bottomrule
    \end{tabularx}
\end{table}

\textcolor{accentColor}{\textbf{Regra:}} Sempre suba criativos na \textbf{sexta-feira} e deixe rodar até \textbf{segunda de manhã}. Muitos criativos campeões só aparecem no 4º dia --- testar por apenas 2 dias é desperdiçar potencial.

\needspace{8cm}

\subsection{Fase 4: Validação e Escala}

Após os testes, você terá dados sobre qual combinação funciona:

\begin{table}[H]
    \centering
    \renewcommand{\arraystretch}{1.4}
    \begin{tabularx}{\textwidth}{X c}
        \toprule
        \rowcolor{tableHeader} \textbf{Resultado do Teste} & \textbf{Ação} \\
        \midrule
        Estrutura ABO melhor na Conta A & Escalar ABO na Conta A \\
        \rowcolor{lightGray}
        Orçamento R\$17 com ROI 5-7x & Duplicar dezenas de campanhas R\$17 \\
        CBO melhor na Conta B & Escalar CBO na Conta B \\
        \rowcolor{lightGray}
        Criativo X com ROAS $\geq$ 1.8 & Escalar com estrutura + conta vencedoras \\
        Nenhum criativo validou & Novos criativos (hooks diferentes) + nova conta \\
        \bottomrule
    \end{tabularx}
\end{table}

\begin{quote}
\textcolor{primaryColor}{\textbf{``Você está diminuindo a quantidade de risco e erro. Testou orçamento, deu essa estrutura, deu bom nessa conta. O que vai fazer? Pegar a estrutura certa, o orçamento certo, na conta certa e rodar. A chance de errar cai drasticamente.''}}
\end{quote}

\needspace{10cm}

\section{Resumo Operacional: Checklist de Implementação}

\subsection{Semana 1: Concepção e Setup}

\begin{itemize}[label=$\square$]
    \item Definir eixo da oferta (timing sazonal ou modelagem de mercado)
    \item Fazer spy e identificar mecanismo validado
    \item Criar aplicativo MVP (Replit/Lovable)
    \item Definir nome chiclete (banco de 10 opções)
    \item Montar quiz de personalização (3 blocos)
    \item Configurar 2 contas de anúncio
    \item Produzir 3-5 criativos (estilo orgânico, até 50s)
\end{itemize}

\subsection{Semana 2: Testes}

\begin{itemize}[label=$\square$]
    \item Sexta: subir todos os testes (estrutura + orçamento + conta)
    \item Segunda: primeira análise --- pausar perdedores claros
    \item Terça-Quarta: confirmação de constância
    \item Identificar: melhor estrutura, melhor orçamento, melhor conta
    \item Produzir novos criativos com hooks diferentes (se necessário)
\end{itemize}

\subsection{Semana 3: Escala Inicial}

\begin{itemize}[label=$\square$]
    \item Escalar criativos validados na estrutura + conta + orçamento vencedores
    \item Testar estratégia R\$17 com duplicação em massa
    \item Testar novos nomes chiclete para atingir novas fatias de público
    \item Monitorar CPA e ROAS diário
\end{itemize}

\subsection{Semana 4+: Escala Consolidada}

\begin{itemize}[label=$\square$]
    \item Manter campanhas lucrativas rodando sem intervenção excessiva
    \item Teste semanal de novos criativos (hooks do TikTok)
    \item Rodar múltiplos nomes chiclete em paralelo
    \item Avaliar criação de novos produtos no mesmo nicho (holding)
    \item Considerar transformar o produto em recorrência/SaaS
\end{itemize}

\pagebreak

\section{Erros Fatais Específicos do Low Ticket}

\subsection{1. Entregar E-book como Produto}

\textbf{Erro:} Modelar uma oferta validada de Nutra/VSL e entregar um e-book.

\textbf{Por que é fatal:} O lead percebe valor baixo. Não há recorrência, não há LTV, não há retenção. O quiz promete personalização e o e-book não entrega.

\textbf{Correção:} Sempre entregar \textbf{aplicativo, plano personalizado ou desafio}. É o formato validado.

\subsection{2. Usar o Mesmo Nome Chiclete da Oferta Modelada}

\textbf{Erro:} Copiar a oferta e manter o nome original.

\textbf{Por que é fatal:} Você compete diretamente com o original. O público já viu. Não há percepção de novidade.

\textbf{Correção:} Mude o nome chiclete. Mesmo mecanismo + nome novo = oferta que parece inovadora.

\subsection{3. Criativo com Pessoas Diferentes}

\textbf{Erro:} Usar uma mulher loira no início do criativo e uma morena no antes/depois final.

\textbf{Por que é fatal:} Quebra a conexão e credibilidade. O lead não associa o resultado à pessoa.

\textbf{Correção:} Sempre a mesma pessoa em todo o criativo. Consistência visual gera confiança.

\subsection{4. Apostar em Uma Única Estrutura/Conta}

\textbf{Erro:} ``O mercado disse que CBO 1-3-3 é o melhor, vou usar só isso.''

\textbf{Por que é fatal:} Cada conta performa diferente. Você pode estar descartando uma oferta lucrativa por testar na estrutura errada.

\textbf{Correção:} Sempre fazer o tripé de testes (estrutura + orçamento + conta).

\subsection{5. Testar Criativos por Apenas 2 Dias}

\textbf{Erro:} Criativo não vendeu em 2 dias, pausar imediatamente.

\textbf{Por que é fatal:} Muitos criativos campeões só performam a partir do 4º dia. Fatias de público ruins nos primeiros dias não significam que o criativo é ruim.

\textbf{Correção:} Suba na sexta, deixe rodar até segunda (mínimo). Ideal: 4+ dias de teste.

\pagebreak

\section{Considerações Finais}

O processo de low ticket não é complexo --- é \textbf{metódico}. A taxa de assertividade de 80\% em novas ofertas vem da disciplina em seguir cada etapa sem pular passos.

\subsection{Os 5 Princípios Inegociáveis}

\begin{enumerate}
    \item \textbf{Personalização é 80/20:} O quiz de personalização supera qualquer outro formato. Faça o lead se enxergar no futuro.

    \item \textbf{Entregável em Aplicativo:} E-book não sustenta operação. App gera recorrência, LTV e valuation.

    \item \textbf{Nome Chiclete Novo:} Mesmo mecanismo validado, nome diferente. Múltiplos nomes = múltiplas fatias de público.

    \item \textbf{Criativo Orgânico de 50s:} TikTok style, antes/depois 3x, mesma pessoa, hook viral, CTA de teste grátis.

    \item \textbf{Tripé de Testes:} Estrutura + orçamento + conta. Nunca aposte em uma única configuração.
\end{enumerate}

\vspace{0.5cm}

\begin{center}
\colorbox{lightGray}{%
    \parbox{0.85\textwidth}{%
        \centering
        \vspace{0.3cm}
        \textcolor{primaryColor}{\textbf{``Depois que você tem uma oferta validada,}}\\
        \textcolor{primaryColor}{\textbf{o jogo é criativo e tráfego.}}\\
        \textcolor{primaryColor}{\textbf{E o segredo do tráfego é testar o que funciona pra SUA conta,}}\\
        \textcolor{primaryColor}{\textbf{não o que funciona pro mercado.''}}\\
        \vspace{0.3cm}
    }%
}
\end{center}

\pagebreak

% --- PÁGINA FINAL ---
\thispagestyle{empty}
\vspace*{0.5cm}

\begin{center}
    {\Huge \textcolor{primaryColor}{\textbf{MAPA DE AÇÃO RÁPIDA}}} \\
    \vspace{1cm}
    \rule{12cm}{1.5pt}
    \vspace{1cm}

    \colorbox{accentColor}{%
        \parbox{0.92\textwidth}{%
            \centering
            \vspace{0.8cm}
            \textcolor{white}{\Huge \textbf{APLICAÇÃO IMEDIATA}} \\
            \vspace{0.3cm}
            \textcolor{white}{\Large 4 Semanas para Oferta Validada e Escalando} \\
            \vspace{0.8cm}
        }%
    }

    \vspace{1cm}

    \begin{flushleft}
    \large
    \textbf{Semana 1 --- Concepção:} \\
    \vspace{0.3cm}
    \hspace{1cm} \textcolor{accentColor}{\Large \textbf{Oferta + App + Quiz}} \\
    \hspace{1cm} \textcolor{secondaryText}{Spy $\rightarrow$ Mecanismo $\rightarrow$ Nome Chiclete $\rightarrow$ App MVP $\rightarrow$ Quiz 3 Blocos} \\

    \vspace{1cm}

    \textbf{Semana 2 --- Criativos + Testes:} \\
    \vspace{0.3cm}
    \hspace{1cm} \textcolor{accentColor}{\Large \textbf{Produção + Tripé de Testes}} \\
    \hspace{1cm} \textcolor{secondaryText}{3-5 criativos orgânicos $\rightarrow$ Subir sexta $\rightarrow$ Testar estrutura/orçamento/conta} \\

    \vspace{1cm}

    \textbf{Semana 3 --- Validação:} \\
    \vspace{0.3cm}
    \hspace{1cm} \textcolor{accentColor}{\Large \textbf{Identificar Vencedores}} \\
    \hspace{1cm} \textcolor{secondaryText}{Melhor estrutura + melhor orçamento + melhor conta = configuração vencedora} \\

    \vspace{1cm}

    \textbf{Semana 4 --- Escala:} \\
    \vspace{0.3cm}
    \hspace{1cm} \textcolor{accentColor}{\Large \textbf{Duplicar e Escalar}} \\
    \hspace{1cm} \textcolor{secondaryText}{Estratégia R\$17 em massa + múltiplos nomes chiclete + novos hooks semanais} \\
    \end{flushleft}

    \vspace{0.5cm}

    \colorbox{lightGray}{%
        \parbox{0.9\textwidth}{%
            \centering
            \vspace{0.4cm}
            \textcolor{primaryColor}{\large \textbf{Processo $>$ Talento}} \\
            \vspace{0.2cm}
            \textcolor{secondaryText}{\normalsize 80\% de assertividade vem de seguir o processo, não de intuição} \\
            \vspace{0.4cm}
        }%
    }

    \vfill

    \textcolor{secondaryText}{\footnotesize Documento compilado em Março de 2026} \\
    \textcolor{secondaryText}{\footnotesize Baseado em metodologia de operação real com 140K+ alunos e escala de 100K+/dia}

\end{center}

\end{document}
